\documentclass[a4paper,12pt]{article}
\usepackage[utf8]{inputenc}
\usepackage[margin=2.5cm]{geometry}
%\usepackage{geometry}[2cm,2cm,2cm,2cm]
\usepackage[utf8]{inputenc}
\usepackage[nosumlimits]{amsmath}
\usepackage{amssymb,stmaryrd,amsthm,MnSymbol}
\usepackage{amsfonts,amssymb,amsopn}
\usepackage[colorlinks=true,
            linkcolor=red,
            filecolor=blue,
            citecolor=blue,
            urlcolor=blue,]{hyperref}
\usepackage{color}
\usepackage[dvipsnames]{xcolor}
\usepackage{comment}
\usepackage{graphicx}
\usepackage{natbib}
\bibliographystyle{elsarticle-harv}
\usepackage{multirow}

\usepackage{soul}
\setstcolor{red}
\newcommand{\werner}[1]{{\color{magenta}WB says: #1}}
\newcommand{\colin}[1]{{\color{blue}CJC says: #1}}

\newcommand{\todo}[1]{\vspace{5 mm}\par \noindent
	\framebox{\begin{minipage}[c]{0.95 \textwidth}
			\tt #1 \end{minipage}}\vspace{5 mm}\par}

\DeclareMathOperator{\diff}{d}
%opening
\title{Response to second review}
\author{Werner Bauer, Joshua Hope-Collins, and Colin J. Cotter}

\begin{document}

\maketitle

We thank the reviewer for their careful reading of our revised manuscript
which is much appreciated, a number of important details were spotted. And
apologies that we did not see these ourselves beforehand!

The questions and comments by the Reviewer are written in Roman font,
followed by our responses in italic fonts. In addition, we highlighted
our changes to the revised manuscript in blue.

\begin{enumerate}
\item Equations (1) and (5) do not use a consistent sign of $b$: $D-b$ in (1) and $D+b$ in (5).\\
  \emph{We have settled on ``+'' for both.}
\item Equations (1) and (5): Please define $\nabla^\perp_h$ and check
  the equations. The reference Gibson et al. (2019) states the
  equation using $\nabla^\perp \cdot u$ instead of $\nabla \cdot
  u^\perp$. In particular, please check whether the sign of the
  relevant term is correct. The reference should also be made more
  concrete by stating the section/equation number.\\ \emph{There was
  an error in (1) which we have corrected, and we have defined
  $\nabla^\perp_h$. The reference has also been updated to indicate the
  chapter sections.}

%   \colin{Just checking signs! When we write $\nabla^\perp \cdot u$ on
%     the sphere, we really mean $k\cdot \nabla\times u$ ie the component
%     of the curl in the $k$ direction. Then integration by parts
%     gives
%     \[
%     \int_e (k\cdot \nabla\times u) w\cdot k\times u \diff x
%     = \int_e (\nabla\times [k(w\cdot k\times u)])\cdot u)  \diff x
%     - \int_{\partial e}(k\cdot n\times u) w\cdot k\times u \diff S.
%     \]
%     Which we compare with
%     \[
%     \int_e (\nabla^\perp\cdot u) w\cdot u^\perp \diff x
%     = -\int_e (\nabla^\perp\cdot (w\cdot u^\perp))\cdot u)  \diff x
%     + \int_{\partial e}(n^\perp\cdot u w\cdot u^\perp \diff S,
%     \]
%     which is what we put in the paper. Hence, we need
%     \[
%     \nabla\times [k(w\cdot k\times u)] = -\nabla^\perp (w\cdot u^\perp),
%     \]
%     which works if we have
%     \[
%     \nabla^\perp \chi = -\nabla\times (k \chi),
%     \]
%     which is what I have now added to the text.
%
%     As an aside to Werner, in our implementation, we had $\nabla^\perp\chi
%     = k\times \nabla\chi$, which is worrying at first till we notice that
%     $\nabla\times k$ is zero in the tangent plane. I will try to remember
%     to fix this in future codes or check mathematically whether this holds
%     when we aren't on the sphere.
%   }
\item Page 3: Shouldn't the upwind side be the one with $u \cdot n > 0$, not $u \cdot n < 0$?\\
  \emph{It should. We have corrected it, and reorganised the text a little.}


\item Page 3, footnote 1: $\phi^+ n^+ \cdot u^+ + \phi^+ n^- \cdot u^-$ should be $\phi^+ n^+ \cdot u^+ + \phi^- n^- \cdot u^-$, I assume.\\
  \emph{It should. We have corrected it.}
\item Page 4: "$\mathrm{d}a$ and $\mathrm{d}b$ are the Gateaux derivatives of $a$ and $c$" should be "$\mathrm{d}a$ and $\mathrm{d}c$ are the Gateaux derivatives of $a$ and $c$", I assume.\\
  \emph{Correct, we have corrected it.}
\item Page 6: $l = 1, \dots, N_l$ looks wrong since the upper index
  must not depend on the running index.\\ \emph{It was a confusing
  notation since $N_l$ was the range for $l$. We have clarified this
  now in the text.}
\item Should $W_l$ be defined using an additional product $\prod_{i=1}^s$ since it couples all stages?\\
  \emph{Yes, we have clarified  this.}
\item Equation (24): Should the argument $w$ of $c_L$ be the test function $\phi$?\\
  \emph{Correct, we have corrected this.}
\item Page 10: Which time integration method is used for the $\Delta t = 1s$ reference solution?\\
\emph{It was an IMEX scheme. We have added this.}
\item Page 11: Since $D$ is discretized in $Q$ ($p = 1$ discontinuous
  space on triangles), I would expect three degrees of freedom per
  triangle, i.e., $60 \times 4^p$ layer depth DOFs on level $p$
  instead of $80 \times 4^p$.\\
  \emph{Correct, we have corrected this.}
%   \colin{Checking the maths. Reviewer is correct on layer depth, it should be
%     $3\times 20$ not $4\times 20$. For velocity, BDM2 has 9 edge DOFs per cell
%     but these are shared between two cells, so we divide by 2. Then there are 3
%     interior values, leading to 4.5+3=7.5 which we multiply by 20 to get 150.}
\item Page 12, ``they have one less stage per order'': Shouldn't this be the other way round, i.e., the (classical) order is one less for the same number of stages?~\\
  \emph{These seem to us to be equivalent, but we have rephrased it in case
  it wasn't clear what ``they'' was.}
\end{enumerate}





\end{document}
